\section{Introduction}

Foundation-model applications are no longer single models queried in isolation. A deployed assistant, analyst tool, or enterprise chatbot may combine retrieval-augmented generation (RAG), dense and lexical indexes, rerankers, prompt templates, supervised fine-tuning, LoRA adapters, -data generation, feedback traces, caches, and human-facing citations. RAG connects language models to external knowledge sources~\cite{lewis2020rag,karpukhin2020dpr,izacard2021fid,gao2024ragsurvey}, while instruction tuning and adapter-based specialization make private or domain-specific records part of model behavior. This improves utility but creates a core trustworthy-AI problem: after a deletion request, an operator may remove the original record while the deployed system still behaves as if the record remains through retrieved chunks, embeddings, summaries, question--answer pairs, cache entries, adapters, or evaluation traces.

Privacy and governance regimes motivate data erasure and deletion accountability~\cite{voigt2017gdpr,cohen2022deletioncontrol,fan2025fedfmchallenges}, and machine unlearning provides the algorithmic foundation for making systems forget selected data~\cite{cao2015machineunlearning,ginart2019making,bourtoule2021machineunlearning,guo2020certified,maini2024tofu,yao2024largeunlearning,liu2024rethinking,wang2024llmunlearning,jin2024surveyunlearning}. But in modern LLM and RAG deployments the deletion problem is not confined to the parameter vector: a system can pass a model-centric unlearning test while still retrieving a stale chunk, citing a shadow copy, answering from a derivative, or paraphrasing a fact learned through an adapter. Extraction and membership-inference work shows that such dependence can be exposed through black-box interaction~\cite{carlini2019secretsharer,carlini2021extracting,carlini2023quantifying,shokri2017membership,yeom2018privacy,mattern2023membership}, and RAG evaluation studies retrieval quality and attribution~\cite{gao2023retrievalaugmented,asai2024selfrag,yu2024ragevalsurvey,akkiraju2024facts}; yet none directly answers the operator's audit question: after a deletion request, what measurable evidence shows that the deleted data and its downstream derivatives no longer influence retrieval, generation, or fine-tuned behavior?

We argue that deletion verification should be treated as \emph{counterfactual behavioral indistinguishability}. The relevant question is not whether a chosen leakage string appears, but whether the post-deletion system behaves like a clean reference system $S_{\mathrm{clean}}$ rebuilt as if the deleted record had never entered any retriever, cache, -data generator, adapter, or evaluation trace. This avoids circularity: diagnostic channels can localize failures, but the gold standard is statistical distance from the clean counterfactual. A practical audit then needs three ingredients: provenance (which artifacts are downstream of the target), efficient adaptive testing over high-risk target--channel arms, and behavioral tests that compare $S_-$ with $S_{\mathrm{clean}}$ and explain whether the residual distance is retrieval-, parametric-, cache-, or derivative-mediated.

This paper proposes \method{} (\emph{Auditable Data-Centric Unlearning}), a framework for deletion-risk auditing in RAG and fine-tuned foundation-model systems. \method{} represents deletion targets and their downstream artifacts in a provenance graph; defines an $(\varepsilon,\delta)$ behavioral-deletion guarantee as bounded distance between $S_-$ and $S_{\mathrm{clean}}$ over an adversary-aware probe class; allocates probes through a graph-structured sequential test; and uses direct, paraphrase, retrieval, counterfactual, extraction, trigger, and graph-pivot probes as diagnostic localizers rather than as the failure definition. Instead of a binary claim of perfect forgetting, it returns a confidence-bounded pass, fail, or escalate decision under an operator-chosen tolerance. The goal is not to replace unlearning algorithms, certified removal, differential privacy, or legal review~\cite{dwork2006dp,guo2020certified,kirchenbauer2023watermark}, but to provide measurable technical evidence about residual behavioral dependence in the pipeline surrounding those mechanisms.

We instantiate the framework in \bench{}, spanning four settings: \emph{RealRAG} (public FEVER retain corpus with injected private targets), \emph{NaturalFEVER} (public FEVER passages as deletion units with redacted outputs), \emph{FineTuneSim} (a controlled adapter-memory simulator), and \emph{HybridReal} (RAG records plus derivatives that survive through summaries, caches, examples, or model-side behavior). We further add LoRA-SFT and cached PEFT/LoRA validations, pretrained LoRA margin analysis, TOFU-style forget/retain evaluation, neural dense retrieval (MiniLM/E5/BGE), clean-counterfactual two-sample tests, and component interventions for causal mediation. Across the repeated harness, \method{} detects configured deletion failures in RealRAG, NaturalFEVER, and HybridReal with zero configured false alarms; in FineTuneSim it detects 72.2\% of all configured failures, or 96.3\% of residual adapter-memory failures once utility-only over-unlearning cases are separated as escalations. The results do not prove legal compliance or absolute absence of influence; they show that deletion audits become substantially more informative when they are clean-counterfactual, provenance-aware, graph-adaptive, and multi-channel.

This paper makes the following contributions:
\begin{itemize}
\item We formulate deletion verification for RAG and fine-tuned foundation-model systems as $(\varepsilon,\delta)$ counterfactual behavioral indistinguishability from a clean reference pipeline.
\item We define deletion risk as clean-counterfactual behavioral distance and use direct, paraphrase, retrieval, extraction, and watermark evidence as diagnostic localizers rather than circular ground truth.
\item We present a graph-structured adaptive audit algorithm with confidence bounds over provenance-linked target--channel arms, plus causal intervention diagnostics that decompose residual risk into retrieval-, cache-, derivative-, and parametric-mediated components.
\item We introduce \bench{}, a reusable benchmark spanning private deletion, natural public-evidence deletion, RAG retention, fine-tuning residues, hybrid derivative retention, dense retrieval, pretrained LoRA, and TOFU-style forget/retain evaluation, with an attack and baseline suite showing where exact-match, canary-only, retriever-only, membership-inference, extraction-only, and influence-proxy audits fail.
\end{itemize}

\noindent Section~2 reviews related work; Section~3 defines the audit problem and threat model; Sections~4--5 present the deletion-risk metric and audit algorithm; Section~6 describes \bench{}; Section~7 reports experiments; Sections~8--9 discuss limitations and reproducibility; Section~10 concludes.
